SQL has become the implementation language of modern data infrastructure. It encodes the reporting, analytics, and transformation logic that organizations run every day, is increasingly generated by business intelligence tools and AI assistants rather than written by hand, and accumulates into recurring pipelines of hundreds or thousands of interdependent queries. In cloud data warehouses that charge for computation, the quality of this SQL translates directly into speed and cost. Yet much of it is poorly written, and a poorly written query can defeat even a sophisticated optimizer: the optimizer selects the best plan for the query it is given, not the best query its author could have written. Recovering that better query, a semantically equivalent form that runs faster, is the task of query rewriting. Today this task falls either to human experts, whose attention cannot cover the sheer volume of machine-generated SQL, or to rule-based rewriters, which recognize only the patterns their rules anticipate and thus miss the unfamiliar queries that need help most.

This dissertation develops techniques that use program synthesis and large language models to rewrite SQL beyond predefined rules, automatically discovering semantically equivalent and substantially faster alternatives for individual queries and, ultimately, for entire SQL pipelines.

One way to move beyond rules is to stop applying transformations and start searching: program synthesis explores the space of semantically equivalent SQL directly, so that in principle any better formulation of a query is discoverable, whether or not anyone has captured it as a rule. SlabCity realizes this idea, introducing query dataflows to steer the search toward promising candidates and layered equivalence checking to admit only correct ones. It establishes that rule-free, whole-query rewriting is practical and that it uncovers optimizations rule-based systems structurally cannot reach.

Search alone, however, falters as queries grow complex, and complex queries are where rewriting matters most. Large language models offer what search lacks: broad knowledge of SQL and of how experienced engineers optimize it, though with no guarantee that their suggestions are correct or actually faster. GenRewrite supplies the missing discipline. It distills each successful rewrite into a Natural Language Rewrite Rule, a human-readable strategy it reuses to guide future queries, and it iteratively corrects the model's mistakes until a rewrite is both equivalent and beneficial. The result is an LLM-based rewriter that grows more capable with every query it optimizes and surpasses both rule-based systems and direct prompting.

Even so, both techniques share a blind spot with every rewriter before them: they optimize one query at a time, while modern analytics runs as recurring pipelines whose queries feed one another. A pipeline's dependencies are not merely execution constraints but optimization signals, revealing redundant work, unused computation, and misplaced effort that no single-query rewriter can see. DAGSmith reads these signals to refactor the pipeline as a whole, verifying every change and weighing it jointly with decisions about what to materialize and how often to refresh, and thereby achieves gains far beyond what optimizing each query in isolation can deliver.

Together, these systems trace a path from synthesis to LLMs, and from individual queries to whole pipelines. Rules capture only the rewrites we already know how to write down; this dissertation shows that the space of equivalent SQL beyond them is rich, reachable, and well worth exploring.



\begin{comment}
    

% Phase-1 draft abstract, adapted from the thesis proposal abstract.
% TODO: revise, and check Rackham's abstract length limit (550 words).
This dissertation studies how program synthesis and large language
models (LLMs) can rewrite SQL beyond existing rule-based optimizers,
progressing from single-query transformations to full SQL pipeline
refactoring.

SlabCity is a whole-query, synthesis-based SQL rewriter that searches
for semantically equivalent, faster queries without relying on
predefined rules. It defines dataflows for SQL queries and uses them
to stratify the synthesis search space via a novel dataflow-based
scoring function that focuses the search on promising rewrites, within
a counterexample-guided inductive synthesis loop that resorts to SMT
verification only when needed. On curated LeetCode and Calcite
workloads, it rewrites more queries and achieves multi-fold latency
reductions compared to rule-based systems.

GenRewrite targets the same single-query setting but replaces
synthesis-based search with LLM-guided generation anchored by
Natural-Language Rewrite Rules (NLR2s), which capture successful
rewrite strategies in human-readable form and transfer rewriting
knowledge across queries. Across TPC-DS and JOB (including their
SQLStorm-generated variants), GenRewrite attains higher coverage and
larger speedups than both rule-based systems and baseline LLM
prompting.

DAGSmith lifts these ideas from single queries to pipelines of
interdependent SQL models. Given a workload represented as a query
DAG, it combines structural graph analysis, LLM-guided refactoring
with data-backed equivalence checking, learned cost modeling,
materialization tuning via iterative ILP, and frequency-aware
scheduling to produce output-equivalent pipelines that are
significantly cheaper and faster end to end.

Together, these systems trace a path from synthesis to LLMs for SQL
query rewriting beyond rules, delivering substantial performance gains
from individual queries to large, dependency-rich SQL pipelines.
\end{comment}